%%% SECTION 4: ADAPTED ICH E6(R3) FOR PHYSICAL AI CLINICAL PRACTICE %%%

The Adaption: ICH Harmonised Guideline \cite{ich-adapt} establishes
comprehensive good clinical practice principles specifically designed for
Physical AI oncology clinical trials, building on the internationally recognized
ICH E6(R3) framework \cite{ich-e6r3}. By adapting an existing standard that is
accepted across dozens of countries rather than creating an entirely new
guideline, this work adds credibility to Physical AI implementations through
association with decades of established GCP practice. The adaptation spans four
major sections (Principles, Investigator Responsibilities, Sponsor
Responsibilities, Data Governance), three appendices (System Documentation,
Clinical Trial Protocol, Essential Records), and a comprehensive Physical AI
glossary with 30 specialized definitions.

\subsection{Principles of Physical AI Clinical Practice}
\label{subsec:ich-principles}

The adapted ICH E6(R3) establishes eight foundational principles that govern
Physical AI clinical practice. These principles extend the traditional GCP
framework by requiring that Physical AI systems meet the same ethical and
quality standards as human-directed clinical practice, with additional
safeguards for robotic safety, AI transparency, and system validation.

The first principle establishes that clinical trials involving Physical AI
systems must be conducted in accordance with ethical principles originating in
the Declaration of Helsinki. Robotic interventions, AI-assisted diagnostics,
and automated treatment delivery must preserve participant autonomy and welfare
above all other considerations. The second principle requires that foreseeable
risks from robotic systems, AI algorithms, and automated processes be weighed
against anticipated benefits before trial initiation. Risks specific to robotic
malfunction, AI prediction error, and cyber-physical system failure must be
explicitly evaluated.

The third principle reinforces that participant rights, safety, and well-being
prevail over scientific and societal interests. This extends to interactions
with all Physical AI system categories: surgical robots (da Vinci, Hugo RAS,
Versius), collaborative robots (Franka Panda, Kinova Gen3, xArm 7), humanoid
robots (Atlas, Optimus, Digit), and all therapeutic, diagnostic, assistive, and
rehabilitative robotic systems. The fourth principle requires specialized
training for all individuals conducting Physical AI trials, supplementing
traditional clinical trial competencies with Physical AI system operation, AI
model interpretation, digital twin management, and simulation validation skills.

The fifth principle requires that trial protocols incorporate specifications
for all robotic systems, AI models, simulation frameworks, and digital twin
configurations, referencing applicable USL scores \cite{usl-standard} and
establishing minimum readiness thresholds. The sixth principle requires IRB or
IEC approval with specific review of Physical AI components including robotic
safety parameters, AI validation evidence, and cybersecurity measures. The
seventh principle establishes that Physical AI systems operate as tools under
human oversight, with continuous monitoring and intervention capability
required for autonomous robotic actions. The eighth principle requires informed
consent that clearly describes robotic interactions, AI-assisted decision
making, and digital twin-based treatment planning.

\clearpage

\begin{table}[H]
\centering
\caption{Adapted ICH E6(R3) Foundational Principles for Physical AI}
\label{tab:ich-principles}
\small
\begin{tabularx}{\textwidth}{c X X}
\toprule
\textbf{No.} & \textbf{Traditional GCP Principle} & \textbf{Physical AI Extension} \\
\midrule
1 & Ethical conduct per Declaration of Helsinki & Extends to robotic and AI-assisted procedures \\
2 & Risk-benefit assessment before initiation & Includes robot malfunction and AI error risks \\
3 & Participant rights prevail & Covers all 10 robot type interactions \\
4 & Qualified personnel for all tasks & Physical AI training requirements added \\
5 & Scientific soundness in protocol & Robot specs, USL scores, AI details required \\
6 & IRB/IEC approval with protocol review & Physical AI-specific review mandated \\
7 & Medical care under qualified physician & Human oversight: autonomous robot actions \\
8 & Freely given informed consent & Robot descriptions, AI transparency required \\
\bottomrule
\end{tabularx}
\end{table}

The adaptation also establishes a Physical AI system classification for clinical
trials, organizing systems into ten categories consistent with the USL scoring
methodology \cite{usl-standard} and patient instructions framework
\cite{patient-instructions-paper}. Each category has specific GCP requirements
proportional to its level of patient interaction and autonomous decision-making.

\subsection{Investigator Responsibilities in Physical AI Trials}
\label{subsec:ich-investigator}

Section 2 of the adapted ICH E6(R3) defines twelve categories of investigator
responsibilities adapted for Physical AI contexts. These responsibilities
address the shift in clinical practice where robots perform functions
traditionally done by human staff, requiring investigators to develop new
competencies in Physical AI oversight.

Investigator qualifications and training (Subsection 2.1) now require
demonstrated competency in Physical AI system oversight. Investigators must
complete Physical AI-specific training covering robot operational parameters,
AI model interpretation, emergency procedures for robot malfunction, and
human-robot interaction protocols. Adequate resources (Subsection 2.2) now
include robot maintenance, calibration equipment, charging infrastructure, and
technical support staff. Medical care provisions (Subsection 2.3) must account
for Physical AI system integration, ensuring that robot-assisted care meets the
same quality standards as human-only care.

Communication with IRB/IEC (Subsection 2.4) requires ongoing reporting of
Physical AI system performance, including robot uptime, adverse events
attributable to Physical AI systems, software updates, and calibration status.
Protocol compliance (Subsection 2.5) now encompasses robot operational
parameters, AI model boundaries, and approved Physical AI configurations. Any
deviation from approved Physical AI configurations constitutes a protocol
deviation requiring documentation and reporting.

Informed consent procedures (Subsection 2.6) require Physical AI-specific
elements as detailed in the adapted 21 CFR Part 50 \cite{cfr50-adapt}. Record
keeping (Subsection 2.7) extends to Physical AI system logs including robot
telemetry, AI decision records, sensor data, and maintenance records. Safety
reporting (Subsection 2.8) now covers Physical AI adverse events including
robot malfunctions, AI decision errors, sensor failures, and cybersecurity
incidents.

Premature termination or suspension (Subsection 2.9) includes criteria specific
to Physical AI system failures that may necessitate trial pause or termination.
Progress reports (Subsection 2.10) must include Physical AI performance metrics.
Compliance with financial and audit requirements (Subsection 2.11) extends to
Physical AI system procurement and maintenance costs. Final reporting
(Subsection 2.12) must document the complete Physical AI trial experience
including system performance, patient interactions, and lessons learned.

\subsection{Sponsor Responsibilities in Physical AI Trials}
\label{subsec:ich-sponsor}

Section 3 of the adapted ICH E6(R3) defines fifteen categories of sponsor
responsibilities adapted for Physical AI. Sponsors bear primary responsibility
for ensuring that Physical AI systems are safe, validated, and properly
integrated into trial operations.

Quality management systems (Subsection 3.1) must encompass Physical AI system
validation, ongoing performance monitoring, and continuous improvement
processes. The quality management framework must address robot-specific quality
indicators including mechanical precision, sensor accuracy, AI model
performance, and system uptime. Regulatory submission requirements (Subsection
3.2) must include Physical AI system descriptions, validation reports, and
safety assessments alongside traditional IND content.

IRB/IEC review of Physical AI systems (Subsection 3.3) requires sponsors to
provide comprehensive Physical AI documentation including system specifications,
risk assessments, validation evidence, and proposed monitoring plans. Trial
design (Subsection 3.4) must incorporate robotic workflows, specifying which
Physical AI systems perform which functions, the degree of autonomy for each
system, and the human oversight structure.

Monitoring Physical AI system performance (Subsection 3.5) requires sponsors
to establish continuous monitoring systems that track robot performance in real
time. Safety reporting for Physical AI adverse events (Subsection 3.6) follows
reporting timelines and severity classifications established in the adapted
21 CFR Part 312 \cite{cfr312-adapt}. Data handling (Subsection 3.7) must
address the unique characteristics of robot-generated records, including high
data volumes, real-time streaming data, and multi-modal data (telemetry, video,
sensor readings, AI decision logs).

The 24-hour simulation from Clinical Trial Site Complete Documentation
\cite{site-docs} provides evidence that these sponsor responsibilities can be
fulfilled at scale. The simulation demonstrated that sponsors can effectively
manage 29 robots serving 168 patients over 24 continuous hours while
maintaining quality management, safety reporting, and data handling standards
consistent with adapted ICH E6(R3) requirements.

\subsection{Data Governance for Physical AI Trials}
\label{subsec:ich-data-governance}

Section 4 of the adapted ICH E6(R3) addresses data governance across three
subsections: blinding safeguards, data lifecycle elements, and computerized
systems requirements.

Blinding safeguards (Subsection 4.1) present unique challenges in Physical AI
trials because robots are physically visible to participants. The adaptation
establishes methods for maintaining treatment blinding when Physical AI systems
are present, including scenarios where sham robot procedures may be appropriate,
where robot configurations differ between treatment arms, and where the presence
or absence of specific robot types may convey treatment assignment information.

Data lifecycle elements (Subsection 4.2) for Physical AI trials encompass robot
telemetry data (position, force, speed, temperature), sensor readings (patient
vital signs, environmental conditions, imaging data), autonomous decision logs
(AI model inputs, outputs, confidence scores, alternative actions considered),
and imaging data (intraoperative video, diagnostic scans, positioning
verification). Each data category has specific requirements for collection
frequency, storage format, retention period, and access controls.

Computerized systems requirements (Subsection 4.3) for Physical AI platforms
must ensure data integrity, audit trails, and regulatory compliance. FHIR
\cite{hl7-fhir}, DICOM \cite{dicom}, and MCP \cite{mcp-protocol} serve as
interoperability standards supporting Physical AI data governance. The Physical
AI National MCP Servers \cite{national-mcp-paper} provide the infrastructure
implementation of these data governance requirements, with five servers handling
governance, data flow, safety, interoperability, and analytics functions
respectively.

\subsection{Physical AI System Documentation}
\label{subsec:ich-documentation}

Appendix A of the adapted ICH E6(R3) establishes documentation requirements for
Physical AI systems in clinical trials. System specifications must include robot
hardware details (manufacturer, model, serial number, software version,
calibration status), AI model details (architecture, training data provenance,
validation results, performance benchmarks), and integration details (network
configuration, data flow architecture, safety interlock connections).

Validation reports must demonstrate that each Physical AI system performs within
specified parameters under expected clinical conditions. Safety assessments must
identify all foreseeable Physical AI-specific risks and document mitigation
strategies. Maintenance logs must record all service activities, calibration
checks, software updates, and hardware replacements. Software version control
must ensure traceability of all code changes affecting Physical AI system
behavior.

\subsection{Clinical Trial Protocol Adaptations}
\label{subsec:ich-protocol}

Appendix B establishes that clinical trial protocols must be adapted to include
Physical AI system specifications identifying each robot by type, manufacturer,
model, USL score, and approved operational parameters. Robot operational
parameters must specify workspace boundaries, force limits, speed limits,
precision requirements, and environmental conditions. Fallback procedures for
robot failures must define escalation pathways from automatic recovery through
manual override to human takeover, with response time requirements for each
escalation level. Criteria for when human staff must override Physical AI
decisions must be explicitly stated, including clinical judgment thresholds,
safety signal triggers, and patient request protocols.

\subsection{Essential Records for Physical AI Trials}
\label{subsec:ich-records}

Appendix C defines additional essential records required beyond standard GCP
documentation. Robot calibration records must document all calibration activities
with date, method, results, and personnel. AI model version histories must trace
every model update with version number, change description, validation results,
and deployment date. Sensor validation logs must confirm that all sensors
providing clinical data meet accuracy and precision specifications. Physical AI
adverse event reports must use standardized reporting formats that capture both
the clinical impact and the technical cause of each event.

\subsection{Physical AI Glossary and Definitions}
\label{subsec:ich-glossary}

The adapted ICH E6(R3) provides 30 Physical AI-specific definitions that
standardize terminology across the entire National Platform. Key definitions
include Physical AI System (an integrated robotic and AI platform designed for
direct patient interaction in clinical trial settings), Physical AI Adverse
Event (any untoward occurrence attributable to Physical AI system operation
during a clinical trial), Physical AI Operator (a qualified individual
responsible for monitoring and controlling Physical AI systems during clinical
procedures), and USL Score (a quantitative measure of robot readiness for
multi-site clinical trial deployment across four dimensions). These definitions
are cross-referenced with the 22 definitions in the adapted 21 CFR Part 312
\cite{cfr312-adapt} and the 17 definitions in the adapted 21 CFR Part 50
\cite{cfr50-adapt} to ensure consistency across all three adapted standards.

\begin{table}[H]
\centering
\caption{Selected Physical AI Definitions from Adapted ICH E6(R3)}
\label{tab:ich-definitions}
\small
\begin{tabularx}{\textwidth}{l X}
\toprule
\textbf{Term} & \textbf{Definition} \\
\midrule
Physical AI System & An integrated robotic and AI platform designed for direct or indirect patient interaction in clinical trial settings \\
Physical AI Adverse Event & Any untoward occurrence attributable to Physical AI system operation during a clinical trial \\
Physical AI Operator & A qualified individual responsible for monitoring and controlling Physical AI systems \\
Autonomous Mode & Operation where AI algorithms direct robotic actions without real-time human input \\
Supervised Mode & Operation where human operators maintain continuous control over robotic actions \\
Physical AI Safety Matrix & A pre-procedure verification checklist for all safety systems \\
USL Score & Quantitative measure of robot readiness across 4 dimensions (1.0-10.0) \\
Phase 0 Validation & Mandatory simulation phase before human subjects enrollment \\
Digital Twin & Computational model replicating a physical patient or system for simulation \\
MCP Conformance & Compliance level with Model Context Protocol communication standards \\
\bottomrule
\end{tabularx}
\end{table}

The standardization of terminology across all three adapted standards ensures
that investigators, sponsors, regulators, and IRB members can communicate
precisely about Physical AI trial operations. The glossary eliminates ambiguity
that could arise from different stakeholders using the same terms with different
meanings. For example, the precise definition of ``autonomous mode'' versus
``supervised mode'' determines the level of informed consent detail required,
the intensity of real-time monitoring, and the qualifications needed for the
Physical AI operator on duty. By establishing these definitions in the adapted
ICH E6(R3) and cross-referencing them in the other two adapted standards, the
National Platform creates a unified vocabulary for the entire Physical AI
oncology trial ecosystem.

\subsection{Quality Management Systems for Physical AI}
\label{subsec:ich-quality-management}

The adapted ICH E6(R3) establishes a quality management framework that extends
traditional GCP quality assurance to address the unique requirements of Physical
AI systems in clinical trials. This framework operates at three levels: sponsor
quality management (ensuring overall trial quality), site quality management
(ensuring local operational quality), and system quality management (ensuring
individual Physical AI system quality).

Sponsor-level quality management requires sponsors to maintain a Physical AI
Quality Management Plan (QMP) that identifies all Physical AI systems in the
trial, establishes quality indicators for each system, defines monitoring
frequencies and methods, and specifies corrective action procedures when quality
indicators deviate from acceptable ranges. The QMP must be approved by the
sponsor's quality assurance unit and reviewed by the IRB as part of the
Physical AI safety plan evaluation.

Site-level quality management requires each trial site to implement the
sponsor's QMP locally, including assigning quality responsibilities to specific
personnel, conducting regular Physical AI system audits per the QMP schedule,
documenting all quality findings and corrective actions, and reporting quality
trends to the sponsor. The PSL framework (Section~\ref{sec:psl-usl}) Dimension
3 (Operational Readiness) encompasses site-level quality management as a
prerequisite for site activation.

System-level quality management operates continuously at each Physical AI
system, encompassing automated self-checks (power-on diagnostics, sensor
calibration verification, software version confirmation), real-time performance
monitoring (throughput, accuracy, error rates, response times), scheduled
maintenance per manufacturer specifications and site SOPs, and calibration
management ensuring all sensors and actuators operate within specified
tolerances.

\begin{table}[H]
\centering
\caption{Quality Management Hierarchy for Physical AI Trials}
\label{tab:ich-quality}
\small
\begin{tabularx}{\textwidth}{l X X}
\toprule
\textbf{Level} & \textbf{Responsibility} & \textbf{Key Activities} \\
\midrule
Sponsor & Overall trial quality management & QMP development, monitoring oversight, corrective action review, regulatory reporting \\
Site & Local operational quality & QMP implementation, Physical AI system audits, personnel training, local corrective actions \\
System & Individual Physical AI quality & Automated self-checks, real-time monitoring, scheduled maintenance, calibration management \\
\bottomrule
\end{tabularx}
\end{table}

\clearpage

\subsection{Investigator Responsibilities Detailed Mapping}
\label{subsec:ich-investigator-detail}

\begin{table}[H]
\centering
\caption{Complete Mapping of Investigator Responsibilities for Physical AI}
\label{tab:ich-investigator-map}
\small
\begin{tabularx}{\textwidth}{c l X}
\toprule
\textbf{No.} & \textbf{Category} & \textbf{Physical AI-Specific Requirements} \\
\midrule
2.1 & Qualifications and training & Physical AI system oversight certification; minimum 40 hours Physical AI training \\
2.2 & Adequate resources & Robot maintenance budget, calibration equipment, technical support contracts \\
2.3 & Medical care & Human oversight for all robot-assisted treatments; physician override capability \\
2.4 & IRB/IEC communication & Quarterly Physical AI performance reports; immediate notification of system failures \\
2.5 & Protocol compliance & Robot operational parameters within approved ranges; AI model versions per protocol \\
2.6 & Informed consent & Physical AI-specific consent elements per adapted 21 CFR Part 50 \\
2.7 & Record keeping & Robot telemetry logs, AI decision records, sensor data, maintenance records \\
2.8 & Safety reporting & Physical AI adverse events per four-category system; cybersecurity incidents \\
2.9 & Premature termination & Physical AI system failure criteria for trial pause or termination \\
2.10 & Progress reports & Physical AI performance metrics in periodic reports \\
2.11 & Financial compliance & Physical AI procurement and maintenance costs documented \\
2.12 & Final reporting & Complete Physical AI experience documented with lessons learned \\
\bottomrule
\end{tabularx}
\end{table}

The twelve responsibility categories create a comprehensive framework that
ensures investigators are accountable for every aspect of Physical AI system
management within their trials. Categories 2.1 and 2.2 address preparation
(qualified personnel and adequate resources). Categories 2.3 through 2.6
address conduct (medical care, communication, compliance, and consent).
Categories 2.7 through 2.10 address documentation and reporting. Categories
2.11 and 2.12 address administrative and final obligations.
